\section{La constante holographique de couplage arithmético-géométrique}
\label{lib01:acoplamiento}

Le registre dodécaphasique relie deux réalisations du même état : la
lecture positionnelle qui intervient dans la clôture arithmétique et la lecture
qui conserve l’incidence, l’orientation et la direction. Le développement suivant
détermine les opérations que permet d’effectuer le registre déjà produit,
l’information retenue par ses mémoires et les erreurs sous lesquelles
ses relations peuvent être retrouvées.
Sa fonction de référence admet une formulation précise : l’orbite de phase
du registre est une base complète de son porteur, et la mémoire isométrique
conserve exactement le Gram de cette base.

Le point de départ demeure la composition
\[
 \APP\longrightarrow\mathrm{TRIT}\longrightarrow\TPK
 \longrightarrow\text{état enrichi}
 \longrightarrow\text{structure discrète conjointe du continu}.
\]
APP apporte les feuillets et leurs résidus et quotients ; TRIT conserve le régime et
l’orientation ; TPK sélectionne, transporte et actualise en conservant la retenue,
la route, la frontière, la survie et la mémoire. Les cinq constructions du
continu appartiennent à cette production conjointe, développée dans le corps
précédent. On travaille ici après celle-ci, sur sa représentation entière
dodécaphasique et ses lecteurs typés. Ni le support linéaire ultérieur ni
l’ensemble des candidats d’un décodeur ne remplacent le domaine des
histoires engendrées.

\subsection{Objet, système de coordonnées et production du registre}
\label{lib01:registro}

La définition~\ref{lib01:definicion} fixe l’objet engendré et ses lecteurs. On développe maintenant son encodage réversible, sa capacité d’identification opératorielle et la stabilité de ses réalisations incidentielles.

La dénomination exprime une fonction constitutive de l’objet, et non une
identification automatique à un coefficient d’interaction physique. La
canonicité se réfère au profil, à l’origine et à l’orientation déclarés. Une
transformation du système de coordonnées et une mise à jour qui incorpore de nouveaux événements
ne sont pas la même opération. Il n’est pas non plus affirmé que le même vecteur numérique
doive être un invariant nécessaire de tout système.

Le calendrier de cent huit positions est constitué de douze fenêtres nonadiques.
L’extension
\[
 0\longrightarrow C_{12}\longrightarrow C_{108}
 \longrightarrow C_9\longrightarrow0
\]
conserve la retenue de phase ; sa non-scission et la distinction entre retour de
phase et retour de l’état sont démontrées dans
\cref{prop:k-extension}. Le registre d’une histoire retient, dans chaque
fenêtre, la sous-route, son orientation, la retenue et le registre des incidences.
L’ordre causal de l’extraction est
\begin{equation}
 x_{\rm term}\longmapsto\mathcal R_{12}(x_{\rm term})
 \longmapsto(b^{90},b^{120},Q_{\rm TPK})
 \longmapsto U_{\rm sgn}\longmapsto K.
 \label{lib01:cadena-registro}
\end{equation}
Les coordonnées régionales et le registre signé sont deux lectures du même terminal ; il n'est pas postulé que la matrice ternaire visible détermine à elle seule la coordonnée signée.

Pour fixer matériellement la dernière composition, les canaux produits sont
\begin{align*}
 b^{90}&=(-495,-116,-684,108,601,-90,
                -173,-88,313,560,-397,461),\\
 b^{120}&=(-425,-281,-481,671,-255,30,
                 475,-543,680,251,6,-128),\qquad Q_{\rm TPK}=6263.
\end{align*}
Soit \(B_r=\sum_{j=0}^3b^{120}_{r+3j}\) et \(d_{r,j}=b^{90}_{r+3j}\), pour \(r=1,2,3\). Le lecteur harmonique donne
\[
 s_1=\frac{Q_{\rm TPK}+2B_1+B_2}{3},\qquad
 s_2=s_1-B_1,\qquad s_3=s_2-B_2,
\]
\[
 U^{(r)}=(s_r,d_{r,1}-d_{r,3},d_{r,0}+d_{r,2},d_{r,0}-d_{r,2}).
\]
Dans cette sortie compatible, \((B_1,B_2,B_3)=(972,-1073,101)\),
\((s_1,s_2,s_3)=(2378,1406,2479)\), et
\begin{align*}
 U^{(1)}&=(2378,-452,-668,-322),\\
 U^{(2)}&=(1406,998,-204,-28),\\
 U^{(3)}&=(2479,-551,-371,-997).
\end{align*}
La division par trois appartient aux trois orbites du lecteur et la division par quatre qui suit à l'inverse de Hadamard ; ce ne sont pas des coefficients ajustés à une sortie physique. Sur les orbites \((1,4,7,10)\), \((2,5,8,11)\), \((3,6,9,12)\), on applique
\[
 H_4=\begin{pmatrix}
 1&1&1&1\\1&1&-1&-1\\1&-1&1&-1\\1&-1&-1&1
 \end{pmatrix},\qquad H_4^2=4I.
\]
Les produits \(H_4U^{(r)}\) sont, respectivement,
\[
 (936,2916,2484,3176),\quad(2172,2636,232,584),
 \quad(560,3296,3656,2404).
\]
Leur divisibilité par quatre vérifie l’appartenance au sous-réseau de
Hadamard. Après division et restitution de l’ordre naturel, on obtient
\begin{equation}
 K=(234,543,140,729,659,824,621,58,914,794,146,601),
 \qquad\mathbf1^*K=6263.
 \label{lib01:K}
\end{equation}
L’inversion et ses congruences, y compris l’indice \(4096\) de l’image
signée, sont prouvées dans \cref{lem:k-hadamard,prop:k-sello}. La
reconstruction est unique parce que \(H_4^{-1}=H_4/4\), et non parce qu’on aurait
recherché un registre proche d’une valeur connue d’alpha.

Pour éviter une collision de symboles, nous écrirons
\(\Delta_j=I-S^j\) pour les extracteurs entiers antécédents, avec
\((Sx)_i=x_{i+1}\) et des indices modulo douze. Alors
\[
 \Delta_3K=b^{90},\qquad \Delta_4K=b^{120},\qquad Q(K)=6263.
\]
L'extracteur \((\Delta_3,\Delta_4,Q)\) est injectif, de facteurs invariants \(1,\ldots,1,12\), d'après \cref{prop:k-transversal}. Ses conditions d'image entière et sa loi d'accumulation restent celles du registre antécédent. Les compléments normalisés de mémoire qui seront notés \(D_j\) sont d'autres opérateurs, dérivés de ces déplacements et munis d'un codomaine métrique déclaré.

\subsection{Codage périodique, clôture et inversion exacte}
\label{lib01:kappa}

Fixons \(B=1000\), douze positions et
\[
 N(k)=\sum_{j=1}^{12}k_jB^{12-j},\qquad d_B=B^{12}-1,
 \qquad \kappa(k)=\frac{N(k)}{d_B}.
\]
Dans le système de coordonnées canonique,
\begin{equation}
 \begin{aligned}
 \kappa_{\rm ag}
 &=\frac{234543140729659824621058914794146601}{10^{36}-1},\\
 \alpha_A&=A_{reg}-\kappa_{\rm ag},\qquad
 A_{reg}=\pi_{\rm HMT}+e_{\rm HMT}-\varphi_{\rm HMT}-4.
 \end{aligned}
 \label{lib01:clausura}
\end{equation}
Les trois coordonnées régionales sont exprimées avant d’agir dans ce lecteur ;
leur origine commune avec la clôture dodécaphasique n’oblige pas à confondre l’ordre
d’évaluation des coordonnées corrélées. L’égalité est un lecteur
ultérieur de l’état engendré. En isoler une variable une fois ses sorties connues est
une opération inverse légitime, distincte de l’utilisation d’une valeur cible pour
produire le registre.

\begin{proposition}[Réversibilité et rayon du codage]
\label{lib01:inversion-kappa}
Dans le domaine \(\{0,\ldots,999\}^{12}\), la fraction \(\kappa(k)\),
avec une base et une longueur déclarées, détermine exactement \(k\). Il en va de même
de la paire exacte \((A_{reg},\alpha_A)\) lorsqu’elle vérifie
\eqref{lib01:clausura}. Si
\[
 |\widetilde\kappa-\kappa(k)|<\frac1{2d_B},
\]
l'arrondi de \(d_B\widetilde\kappa\) à l'entier le plus proche récupère \(N(k)\) sans erreur. Si les deux termes de la clôture sont approchés séparément, il suffit que la somme de leurs erreurs absolues soit inférieure à \(1/(2d_B)\).
\end{proposition}
\begin{proof}
La multiplication par \(d_B\) retrouve l’entier \(N(k)\).
Douze divisions euclidiennes par \(B\), en conservant les restes et le nombre
de positions, retrouvent ses blocs, y compris les blocs initiaux nuls.
L’hypothèse d’erreur situe l’entier vrai à une distance strictement
inférieure à une demi-unité. Pour la clôture, on applique l’inégalité
triangulaire à l’erreur de la différence. Deux entiers consécutifs montrent
que le rayon ne peut pas être augmenté uniformément sur tout le domaine du
code : leur milieu est ambigu.
\end{proof}

Si l’on reçoit une fraction réduite \(p/q\), on calcule
\(N=d_Bp/q\) et on vérifie l’intégralité et \(0\leq N\leq d_B\).
Le dénominateur réduit ne remplace pas les métadonnées de base, de longueur et de
système de coordonnées. Le code contient \(10^{36}\) mots et nécessite
\(\lceil\log_2(10^{36})\rceil=120\) bits dans une représentation binaire
de longueur fixe. Douze blocs de dix bits occupent également 120 bits : aucune
compression supplémentaire, authentification ni intégrité
cryptographique n’a été démontrée. Une fraction exacte ou l’entier \(N\) évite de perdre la
précision nécessaire à cette inversion.

La lecture périodique ne se confond pas avec la troncature \(N/B^{12}\). Si \(\kappa_j=\kappa(S^jK)\), une division positionnelle donne
\begin{equation}
 \kappa_{j+1}=1000\kappa_j-K_{j+1},\qquad
 K_{j+1}=1000\kappa_j-\kappa_{j+1},\qquad
 \sum_{j=0}^{11}\kappa_j=\frac{6263}{999}.
 \label{lib01:orbita-kappa}
\end{equation}
Ainsi, l’orbite scalaire conserve toutes les différences de phase et permet
de lire \(B_K=\{j+1:1000\kappa_j-\kappa_{j+1}\geq729\}\).
La réversibilité de ces représentations n’implique pas qu’un scalaire
reconstruise toute l’histoire enrichie qui a produit le registre.

\subsection{L’orbite complète comme référence de son porteur}
\label{lib01:orbita}

Dans \(V=\mathbb R^{12}\), muni de son produit euclidien canonique, considérons
\[
 O_K=(K,SK,\ldots,S^{11}K):\mathbb R^{12}\longrightarrow V.
\]
Le décalage \(S\) agit sur les positions du registre. Il n’est pas
identifié par la seule périodicité à l’évolution complète TPK ou à
une symétrie exceptionnelle ultérieure.

\begin{theorem}[Référence cyclique complète et matrice de Gram exacte]
\label{lib01:marco-ciclico}
La matrice \(O_K\) est inversible. Sa matrice de Gram a pour valeurs propres
\[
\begin{array}{c|c}
 \text{valeur}&\text{multiplicité}\\\hline
 39225169=6263^2&1\\
 1248025-345420\sqrt3&2\\
 589609&2\\
 1407529&2\\
 1053157&2\\
 1248025+345420\sqrt3&2\\
 697225&1
\end{array}
\]
et satisfait
\begin{equation}
 589609\|c\|^2\leq\|O_Kc\|^2\leq39225169\|c\|^2,
 \quad
 \|O_K^{-1}\|=\frac1{\sqrt{589609}},\quad
 \operatorname{cond}(O_K)=\frac{6263}{\sqrt{589609}}.
 \label{lib01:cotas-orbita}
\end{equation}
L’orbite du registre centré a rang onze.
\end{theorem}
\begin{proof}
Les entrées de la matrice de Gram sont les autocorrélations \(\langle K,S^jK\rangle\). Dans l'ordre \(j=0,\ldots,11\), elles valent
\begin{align*}
 (&4251257,2999323,3163385,3287920,3216553,3344743,\\
  &2950064,3344743,3216553,3287920,3163385,2999323).
\end{align*}
Comme \(S\) est orthogonal, le gramien est circulant. Substituer les douze
caractères \(\exp(2\pi i m/12)\) dans cette ligne donne le tableau réel
précédent. Toutes les valeurs sont positives. La plus petite des deux valeurs
avec \(\sqrt3\) est supérieure à \(589609\), car
\(658416>345420\sqrt3\), inégalité qui se vérifie en élevant au
carré ses deux membres positifs. Les autres valeurs du tableau sont
supérieures ou égales à \(589609\). Le mode uniforme a pour valeur
\((\sum K_i)^2=6263^2\) ; l’inégalité triangulaire de la transformée
de Fourier, en utilisant \(K_i\geq0\), borne toutes les autres par cette même
valeur. Le théorème spectral donne \eqref{lib01:cotas-orbita}. Centrer
le registre élimine seulement le mode uniforme ; les onze autres conservent
les valeurs non nulles du tableau.
\end{proof}

Le déterminant, avec l’ordre des colonnes et l’orientation de \(S\)
fixés ici, est
\[
 \det O_K=5483119259214020183850397930008283625.
\]
L'absence de période propre du vecteur ne prouverait pas à elle seule l'inversibilité : le point décisif est qu'aucun de ses modes de Fourier ne s'annule. Le conditionnement numérique est approximativement \(8.15643\). Si l'on normalise \(K\) à une norme égale à un, les extrêmes de la matrice de Gram sont divisés par \(\|K\|^2=4251257\) ; l'amplitude de ses composantes ne doit pas être confondue avec une amélioration intrinsèque par rapport à toute référence de même norme.

\begin{corollary}[Séparation polaire de la forme et de l’orientation]
\label{lib01:polar}
Soient
\[
 G_K=O_K^*O_K,\qquad P_K=G_K^{1/2},\qquad
 F_K=O_KG_K^{-1/2}.
\]
Alors \(P_K\) est défini positif, \(F_K\) est orthogonal et
\(O_K=F_KP_K\). Une réalisation isométrique \(J:V\to\mathcal Y\)
conserve \(G_K\) et \(P_K\), et transporte la partie isométrique vers
\(JF_K\).
\end{corollary}
\begin{proof}
La positivité du Gram permet de définir sa racine et son inverse au moyen du
théorème spectral. L’identité
\(F_K^*F_K=G_K^{-1/2}G_KG_K^{-1/2}=I\) prouve l’assertion.
De plus, \((JO_K)^*(JO_K)=O_K^*O_K\).
\end{proof}

Cette décomposition polaire est définie à partir du Gram complet : elle ne s’obtient pas uniquement à partir de la
liste de ses valeurs propres ou de la norme de \(K\). Le symbole
\(P_K\) de ce corollaire n’est pas le projecteur dimensionnel de rang dix
qui sera construit plus bas. Cette séparation est une conséquence de la
référence cyclique, et non une identification de tous ses lecteurs.

\subsection{Identification d'opérateurs et holonomie relative}
\label{lib01:identificacion}

\begin{theorem}[Douze réponses générales ou une réponse cyclique]
\label{lib01:operadores}
Pour tout opérateur linéaire \(H:V\to V\), ses douze réponses \(Y=(HK,HSK,\ldots,HS^{11}K)\) déterminent
\[
 H=YO_K^{-1}.
\]
Si \(HS=SH\), une unique réponse vectorielle complète \(y=HK\) détermine
\begin{equation}
 H=O_yO_K^{-1},\qquad
 H=\sum_{j=0}^{11}h_jS^j,\qquad h=O_K^{-1}y.
 \label{lib01:filtro}
\end{equation}
Avec des erreurs \(E\) dans la matrice des réponses ou \(\delta y\) dans le
vecteur, respectivement,
\begin{align}
 \|\widehat H-H\|&\leq\frac{\|E\|}{\sqrt{589609}},
 \label{lib01:error-matriz}\\
 \|\widehat h-h\|&\leq\frac{\|\delta y\|}{\sqrt{589609}},\qquad
 \|\widehat H-H\|\leq\sqrt{\frac{12}{589609}}\,\|\delta y\|.
 \label{lib01:error-filtro}
\end{align}
\end{theorem}
\begin{proof}
La première égalité est \(Y=HO_K\). Si \(H\) commute avec \(S\), alors \(HS^jK=S^jy\), de sorte que \(Y=O_y\). La commutation avec le cycle détermine toutes les colonnes de \(H\) à partir de l'une d'entre elles ; les douze puissances \(I,S,\ldots,S^{11}\) sont indépendantes et engendrent cet espace d'opérateurs. Par conséquent, \(y=O_Kh\), ce qui donne la seconde formule. Les deux premières bornes utilisent \(\|O_K^{-1}\|=1/\sqrt{589609}\). La dernière utilise \(\|O_{\delta y}\|\leq\|O_{\delta y}\|_{\rm F}
=\sqrt{12}\|\delta y\|\).
\end{proof}

Une réponse signifie ici un vecteur à douze composantes, non une
mesure scalaire. Par exemple, si \(H=2I+3S-S^5\), la réponse
\(y=HK\) restitue exactement les coefficients \(2,3,-1\) à leurs
positions au moyen de \eqref{lib01:filtro} ; l’inverseur a besoin de \(y,K,S\),
non des coefficients du filtre. En revanche, l’opérateur non nul dont la première
ligne est \((543,-234,0,\ldots,0)\) et dont les autres lignes sont nulles
annule \(K\). Une seule réponse n’identifie pas un opérateur arbitraire.

La fonction de référence admet une généralisation exacte. Dans un cycle de
\(n\) positions, \(H\mapsto Hg\) identifie tous les opérateurs
qui commutent avec le cycle si et seulement si
\(O_g=(g,Sg,\ldots,S^{n-1}g)\) est inversible. La suffisance est la
preuve précédente. Si \(O_g\) est singulier, un vecteur non nul de son noyau
donne un polynôme de degré inférieur à \(n\) qui annule \(g\) ; l’opérateur
est non nul car les puissances du cycle sont indépendantes. L’impulsion
unitaire a également une orbite orthonormale et sert de référence. La
conséquence spécifique pour \(K\) est que le même objet déjà utilisé dans
la clôture et l’incidence vérifie cette condition, sans avoir été choisi de
nouveau pour optimiser l’essai. On n’en déduit pas une identification de
dynamiques non linéaires ou d’états cachés arbitraires ; une application à
une linéarisation exige que celle-ci ait été définie et justifiée.

\begin{corollary}[Récupération de l'holonomie relative à partir du complément]
\label{lib01:holonomia}
Soient \(H_B\) connu et inversible et \(H_C\) un autre transport du même type. Le complément
\[
 D_\gamma=\frac{\sqrt8}{9}(H_C-H_B)
\]
détermine
\[
 H_B^{-1}H_C=I+\frac9{\sqrt8}H_B^{-1}D_\gamma.
\]
Si \(D_\gamma\) commute avec \(S\), en particulier si les deux
transports le font, \(y=D_\gamma K\) suffit à récupérer
\begin{equation}
 H_B^{-1}H_C
 =I+\frac9{\sqrt8}H_B^{-1}O_yO_K^{-1}.
 \label{lib01:holonomia-una-respuesta}
\end{equation}
L'erreur dans cette reconstruction est bornée par
\[
 \frac9{\sqrt8}\|H_B^{-1}\|
 \sqrt{\frac{12}{589609}}\,\|\delta y\|.
\]
\end{corollary}
\begin{proof}
On isole \(H_C\) dans la définition du complément et l'on applique \cref{lib01:operadores}. La borne résulte de la sous-multiplicativité.
\end{proof}

Pour \(H_B=S^3\) et \(H_C=S^4\), la donnée rationnelle
\(y/\sqrt8=(S^4-S^3)K/9\) permet de reconstruire l’holonomie relative
\(S\) sans fournir \(H_C\) à l’inverseur. Si \(H_B\) est
unitaire, le facteur d’erreur est d’environ \(0.01435510\). Sans
commutation, on maintient les douze réponses générales, non une hypothèse
fictive de symétrie. La sélection d’une boucle physique et sa représentation
appartiennent toujours à leur dynamique génératrice ; ce protocole ne les
remplace pas.

\subsection{La mémoire complète conserve la référence et les observables}
\label{lib01:memoria-completa}

Soient \(C_n\) des opérateurs unitaires d'un espace de Hilbert \(H\),
\[
 T_n=\frac{8I+C_n}{9},\qquad
 D_n=\frac{\sqrt8(C_n-I)}9,\qquad
 R_0=I,\quad R_{n+1}=T_nR_n.
\]
L'écriture inclut des suites non stationnaires. Elle n'exige pas de commutation entre les \(C_n\) d'étapes distinctes.

\begin{theorem}[Registre isométrique, terminal et transport]
\label{lib01:isometria}
Pour tout \(N\),
\[
 R_N^*R_N+\sum_{n=0}^{N-1}R_n^*D_n^*D_nR_n=I.
\]
Il existe la limite forte positive \(A_\infty=\lim_N R_N^*R_N\), et
\begin{equation}
 \mathcal Vx=
 \bigl(A_\infty^{1/2}x,(D_nR_nx)_{n\geq0}\bigr)
 \quad\text{satisfait}\quad\mathcal V^*\mathcal V=I.
 \label{lib01:V-infinita}
\end{equation}
Son image \(\mathcal M\) est fermée et l'inverse sur celle-ci est \(\mathcal V^*\), de norme un. Dans le cas fini, on remplace la première composante par \(R_Nx\) et l'on conserve les \(N\) compléments. Pour tout opérateur borné \(A\),
\[
 A^{\mathcal V}=\mathcal VA\mathcal V^*\big|_{\mathcal M}
\]
préserve les produits, les adjoints, la norme et le spectre dans cette image, et
\[
 \langle\mathcal Vx,A^{\mathcal V}\mathcal Vy\rangle
 =\langle x,Ay\rangle.
\]
\end{theorem}
\begin{proof}
Développer les deux termes d’une étape donne
\(T_n^*T_n+D_n^*D_n=I\). Par conséquent,
\(R_n^*R_n-R_{n+1}^*R_{n+1}=R_n^*D_n^*D_nR_n\), et la somme
télescopique démontre l’identité finie. La suite d’opérateurs
positifs \(R_N^*R_N\) est décroissante et minorée ; sa
limite forte existe. En prenant des produits avec \(x\), la somme des
carrés des normes des compléments est
\(\|x\|^2-\langle x,A_\infty x\rangle\). Cela prouve que la
série appartient à la somme hilbertienne et que \(\mathcal V\) est une
isométrie. Une isométrie a une image fermée et un inverse adjoint sur celle-ci.
Insérer \(\mathcal V^*\mathcal V=I\) dans les compositions démontre
les propriétés du transport. On n’a supposé ni la convergence de
\(R_Nx\), ni que \(A_\infty\) soit un projecteur.
\end{proof}

En appliquant ce résultat au porteur du registre,
\begin{equation}
 (\mathcal VO_K)^*(\mathcal VO_K)=O_K^*O_K.
 \label{lib01:gram-memoria}
\end{equation}
Les douze directions et toutes leurs corrélations sont conservées, avec les
mêmes bornes de \cref{lib01:marco-ciclico}, à toute profondeur de
mémoire complète. On n’affirme pas que le continu complet soit de dimension
douze : cette référence de douze dimensions est conservée dans son image.
Les opérateurs corrects y sont
\(S^{\mathcal V}=\mathcal VS\mathcal V^*\) et
\(H^{\mathcal V}=\mathcal VH\mathcal V^*\). On ne présume pas qu’un
décalage ambiant des cases de stockage soit
\(S^{\mathcal V}\).

Si \(z=\mathcal Vy+\delta z\), l'estimation
\[
 \widehat h=O_K^{-1}\mathcal V^*z
\]
maintient la première borne de \eqref{lib01:error-filtro} avec \(\|\delta z\|\) à la place de \(\|\delta y\|\). La profondeur ne dégrade pas cette constante car l'isométrie et son terminal sont conservés. Une sélection partielle de canaux possède une autre matrice de Gram et requiert ses propres bornes. Les puissances extérieures \(\bigwedge^p\mathcal V\) conservent les déterminants de Gram ; les traces de rang fini et les mesures spectrales des observables transportés sont également conservées. Les carrés des normes ne sont pas pour autant automatiquement interprétés comme des probabilités de Born : une telle lecture exige leur réalisation ultérieure.

\subsection{Incidence exceptionnelle et direction sélectionnée}
\label{lib01:direccion}

Soit \(e=\mathbf1\), \(\Pi=I-ee^*/12=P_{11}\). Dans le système de coordonnées
icosaédrique \(A_5/C_5\) déjà construit, la représentation orthogonale
\(\rho\) et le caractère tridimensionnel fixent
\begin{equation}
 P_3=\frac3{60}\sum_{g\in A_5}\chi_3(g^{-1})\rho(g),
 \quad P_3^*=P_3=P_3^2,\quad
 \operatorname{rank}P_3=3,\quad P_3e=0.
 \label{lib01:P3}
\end{equation}
Ce projecteur est reçu avec sa représentation et ses marques ; il n’est pas choisi
pour adapter une direction au registre. Posons
\[
 k=\Pi K,\qquad u=P_3k,\qquad
 E_u=\frac{uu^*}{\|u\|^2},\qquad Q_K=\Pi-E_u=P_{10}(K).
\]
Les évaluations exactes sont
\begin{align}
 \|k\|^2&=\frac{11789915}{12},&
 \|u\|^2&=\frac{6638585+2275584\sqrt5}{20},
 \label{lib01:normas-direccion}\\
 \eta_K:=\frac{\|u\|^2}{\|k\|^2}
 &=\frac{3(6638585+2275584\sqrt5)}{5\cdot11789915}
 \simeq0.5967954228259091.
 \label{lib01:eta}
\end{align}
En particulier, \(u\ne0\). Comme \(u^*k=\|u\|^2\),
\begin{equation}
 E_uk=u,\qquad Q_Kk=k-u,\qquad
 k=u+Q_Kk,\qquad \|k\|^2=\|u\|^2+\|Q_Kk\|^2.
 \label{lib01:particion-direccion}
\end{equation}
La réduction \(12\to11\to10\) élimine d’abord la direction uniforme
puis une droite sélectionnée. Elle n’élimine pas tout \(\operatorname{im}P_3\) :
dans \(\operatorname{im}Q_K\) subsistent les deux directions de ce
sous-espace orthogonales à \(u\). Ces dimensions sont celles des
espaces algébriques écrits ; leur identification à des dimensions physiques
exige leurs applications de réalisation.

\begin{proposition}[Transport exact de la direction et du descripteur]
\label{lib01:transporte-direccion}
Pour une isométrie complète \(\mathcal V\), soient
\(\Pi^{\mathcal V}=\mathcal V\Pi\mathcal V^*\),
\(P_3^{\mathcal V}=\mathcal VP_3\mathcal V^*\) et
\(u^{\mathcal V}=\mathcal Vu\). Dans \(\mathcal M\),
\[
 Q_K^{\mathcal V}
 =\mathcal VQ_K\mathcal V^*
 =\Pi^{\mathcal V}
 -\frac{u^{\mathcal V}(u^{\mathcal V})^*}{\|u^{\mathcal V}\|^2},
 \quad \operatorname{rank}Q_K^{\mathcal V}=10,
 \quad \eta_K^{\mathcal V}=\eta_K.
\]
En outre, \(\mathcal V^*u^{\mathcal V}=u\) et \(\mathcal V^*Q_K^{\mathcal V}\mathcal V=Q_K\).
\end{proposition}
\begin{proof}
On substitue les définitions et \(\mathcal V^*\mathcal V=I\). Les normes et le rang sont conservés par l'isométrie. L'identité de l'image est \(\mathcal V\mathcal V^*\), et non l'identité d'un espace de stockage ambiant plus grand. Prolonger par zéro en dehors de l'image n'ajoute pas de directions physiques.
\end{proof}

L’écran incidentiel complet offre une autre réalisation exacte. Si
\(\mathcal I\) est l’opérateur d’incidence déjà construit et
\(E_{\rm inc}=\mathcal I(e^\perp)\), alors
\begin{equation}
 F:V\longrightarrow Y=\mathbb R\oplus E_{\rm inc},\quad
 Fx=\left(\mu(x),\frac{\mathcal I\Pi x}{6}\right),\quad
 \mu(x)=\frac{\mathbf1^*x}{12},\qquad
 \langle(\mu,y),(\nu,z)\rangle_Y=12\mu\nu+\langle y,z\rangle
 \label{lib01:incidencia-F}
\end{equation}
est une isométrie surjective, puisque
\(\Pi\mathcal I^*\mathcal I\Pi=36\Pi\). Pour l’incidence
brute des 132 hexades, l’identité complète est
\(\mathcal I^*\mathcal I=36I+30\mathbf1\mathbf1^*\) ; le
centrage élimine exactement sa composante uniforme. L’image incidentielle a
dimension onze ; \(Y\) n’est pas tout \(\mathbb R\oplus\mathbb R^{132}\).
Appliquer \(F\) terme à terme à la mémoire conserve tant le gramien
que les opérateurs précédents et leurs domaines. Cette
composition n’introduit pas d’action du Monstre sur \(K\).

\subsection{Stabilité continue et nécessité de conserver les corrélations}
\label{lib01:ruido}

\begin{theorem}[Erreur de la direction à partir de la mémoire complète]
\label{lib01:error-proyectores}
Soit \(z=\mathcal VK+\delta\), \(\|\delta\|\leq\varepsilon\), et reconstruisons \(\widehat K=\mathcal V^*z\), \(\widehat k=\Pi\widehat K\), \(\widehat u=P_3\widehat k\). Les trois erreurs sont inférieures ou égales à \(\varepsilon\). Si \(\varepsilon<\|u\|\),
\[
 \|Q_{\widehat K}-Q_K\|\leq\frac\varepsilon{\|u\|}.
\]
Si \(\varepsilon<\|k\|\),
\[
 |\eta_{\widehat K}-\eta_K|\leq\frac\varepsilon{\|k\|}.
\]
\end{theorem}
\begin{proof}
L’adjoint de l’isométrie et les projecteurs sont de norme un. Le bruit
orthogonal à \(\mathcal M\) disparaît lorsqu’on applique \(\mathcal V^*\).
Pour deux droites non nulles, la norme de la différence de leurs projecteurs est
le sinus de l’angle entre elles. La distance de \(u\) à la droite de
\(\widehat u\) ne dépasse pas \(\|u-\widehat u\|\), de sorte que ce
sinus est inférieur ou égal à \(\varepsilon/\|u\|\). L’hypothèse
exclut \(\widehat u=0\).

Pour la seconde affirmation, on normalise \(k\) et \(\widehat k\). La différence de leurs projecteurs de rang un a, dans leur plan, les valeurs propres \(\pm\sin\theta\). En l'appariant avec \(0\leq P_3\leq I\), la valeur absolue ne dépasse pas \(\sin\theta\) : dans une base de vecteurs propres, on obtient \(\sin\theta(a-b)\), avec \(a,b\in[0,1]\). Le même argument de distance donne \(\sin\theta\leq\varepsilon/\|k\|\). Si les droites coïncident, la différence des quotients est nulle.
\end{proof}

Ici \(\|u\|\simeq765.733\), de sorte que le facteur pour le projecteur
est d’environ \(0.00130594\). Dans le lecteur à deux mémoires seulement,
l’inverse a pour norme \(9/4\) et le facteur correspondant est
\(9/(4\|u\|)\simeq0.00293836\). Ce sont des données observées distinctes ;
la différence n’est pas une amélioration gratuite du même canal. La correction
entière de \cref{lib02:radios-memoria} permettra de récupérer d’abord le
registre sans erreur à l’intérieur de son rayon, au lieu de propager une erreur
continue jusqu’au projecteur.

\begin{proposition}[Les blocs croisés font partie de l’observable]
\label{lib01:correlaciones}
Écrivons \(\mathcal Vx=(W_ax)_a\), en incluant le terminal. Alors
le bloc \((a,b)\) de \(A^{\mathcal V}\) est
\[
 (A^{\mathcal V})_{ab}=W_aAW_b^*.
\]
Sa forme quadratique se calcule comme limite de troncatures finies
\[
 \langle y,A^{\mathcal V}y\rangle
 =\lim_F\sum_{a,b\in F}\langle y_a,W_aAW_b^*y_b\rangle.
\]
La conservation générale ne s’obtient pas en écartant tous les termes
\(a\ne b\).
\end{proposition}
\begin{proof}
La formule par blocs résulte de la multiplication de \(\mathcal VA\mathcal V^*\). Les projecteurs de troncature convergent fortement vers l'identité et \(A^{\mathcal V}\) est borné, ce qui prouve la limite sans affirmer la convergence absolue d'une série double arbitraire. Pour un contre-exemple, prenons \(H=\mathbb R\), \(C=-I\), \(T=7/9\), \(D=-2\sqrt8/9\), \(\mathcal Vx=(Tx,Dx)\) et \(A=I\). La partie diagonale, évaluée en \(\mathcal Vx\), vaut
\[
 (T^4+D^4)x^2=\frac{3425}{6561}x^2.
\]
Les termes croisés apportent \(3136x^2/6561\), et la somme est
\(x^2\), comme l’exige l’isométrie.
\end{proof}

\subsection{Correction discrète avant la sélection incidencielle}
\label{lib01:correccion-incidencia}

Les lecteurs canoniques à deux mémoires sont
\[
 T_j=\frac{8I+S^j}{9},\qquad
 D_j=\frac{\sqrt8(S^j-I)}9,\qquad
 Mx=(D_3x,D_4T_3x).
\]
La seconde mémoire agit après la première compression ; on ne réinjecte pas
dans ce protocole le premier complément. Le théorème général de
\cref{lib02:radios-memoria} démontrera, sans utiliser le registre comme
cible, les rayons stricts
\[
 r_0=\frac{2\sqrt{146}}{81},\qquad
 r_q=\frac{4\sqrt{73}}{81}.
\]
Le premier rayon retrouve \(\Pi K\) à partir de \(MK\), modulo la
translation entière uniforme. Le second retrouve \(K\) intégralement lorsqu’on
conserve également \(q=\mathbf1^*K\). Si l’on observe une charge réelle
avec une erreur inférieure à \(1/2\), son arrondi retrouve d’abord le canal entier
de charge. Aucun de ces procédés n’a besoin du terminal
\(T_4T_3K\).

La sélection du bloc haut dans le système de coordonnées marqué est
\begin{equation}
 B_K=\{i:K_i\geq729\}=\{4,6,9,10\}.
 \label{lib01:umbral}
\end{equation}
La quatrième composante est exactement au seuil. Cette circonstance
sépare la stabilité continue de la stabilité discrète.

\begin{proposition}[Marge continue nulle et restitution entière]
\label{lib01:umbral-estable}
Le sélecteur direct de \eqref{lib01:umbral}, appliqué à la reconstruction réelle des mémoires, n'a pas de rayon euclidien positif de stabilité en \(MK\), même avec une charge exacte. En revanche, le décodage entier préalable avec une erreur inférieure à \(r_q\) restitue exactement \(B_K\) sans modifier le seuil.
\end{proposition}
\begin{proof}
Soit \(L\) l’inverse des moindres carrés, qui satisfait
\(LM=\Pi\). Pour \(t>0\), posons
\[
 e_t=-tM\Pi e_4,
 \qquad \widetilde K_t=\frac q{12}\mathbf1+L(MK+e_t)
 =K-t\Pi e_4.
\]
On a \(\|e_t\|\to0\), mais la quatrième composante est
\(729-11t/12<729\). Pour \(t\) petit, seule cette appartenance change.
La seconde assertion découle de la récupération exacte préalable de \(K\)
au moyen de \cref{lib02:radios-memoria}, puis de l’application du même sélecteur.
\end{proof}

Si les autres données marquées \(P_\pi,H_e,H_\varphi,H_{\rm APP}\)
sont conservées exactes ou transportées selon leur loi, l’unicité de la
construction exceptionnelle restitue le même drapeau
\(B_K\subset H^-_\alpha\) et la même origine marquée. Corriger le
registre n’équivaut pas à corriger des erreurs indépendantes de ces marques.
Le témoin
\[
 K'=K-e_4+e_6,
 \qquad Q(K')=6263,\qquad B_{K'}=\{6,9,10\}
\]
appartient également à \(\{0,\ldots,999\}^{12}\) et satisfait \(\|MK'-MK\|^2=4672/6561\). Le point milieu des deux mémoires se trouve à distance \(r_q\) de chacune. Il prouve que le rayon strict est optimal uniformément dans le domaine algébrique de charge fixe, et non que \(K'\) possède une histoire terminale HMT construite.

Sans charge, \(M(K-\mathbf1)=MK\) et les sélections hautes sont distinctes.
Les deux mémoires peuvent retrouver exactement \(Q_K\) et \(\eta_K\),
qui ignorent la moyenne, sans retrouver ce sélecteur absolu. Pour le
projecteur, la condition \(P_3\Pi K\ne0\) reste nécessaire ;
la correction entière ne supprime pas cette condition de définition. Un essai
fini concret utilise \(z=MK/\sqrt8\) et ajoute \(1/10\) à sa première
composante. L’erreur dans la norme de \(M\) est \(\sqrt8/10<r_q\).
L’algorithme de plancher et de plafond de \cref{lib02:decodificador} laisse 924
candidats de charge 6263 et retrouve \(K\) et \(B_K\) exactement.

En appliquant l’isométrie \eqref{lib01:incidencia-F}, on obtient
\(M_{\rm inc}=(F\oplus F)M\). Toutes les distances entre mots
et les normes de bruit sont conservées, de sorte que les mêmes rayons et
décodeurs valent dans \(Y\oplus Y\). C’est la connexion
quantitative entre mémoire et incidence ; elle ne dépend pas de l’attribution au
déplacement \(S\) d’une symétrie exceptionnelle supplémentaire.

\subsection{Covariance, raffinement et charge variationnelle}
\label{lib01:covarianza}

Si \(G\) est orthogonal et que l'on transporte conjointement \(K'=GK\), \(\Pi'=G\Pi G^*\) et \(P_3'=GP_3G^*\), alors
\[
 u'=Gu,\qquad Q_{K'}'=GQ_KG^*,\qquad \eta_{K'}'=\eta_K.
\]
Dans un système de coordonnées fixe, \(K\mapsto aK+b\mathbf1\), \(a\ne0\), laisse
invariants \(Q_K\) et \(\eta_K\) ; il ne laisse pas invariant l’objet entier.
Par exemple, \(\kappa(K+\mathbf1)-\kappa(K)=1/999\), et la lecture
\(\alpha_A\) changerait de \(-1/999\) si les
régions étaient maintenues fixes. C’est un contre-exemple algébrique à l’identification du descripteur à
l’ensemble du registre, et non une affirmation d’admissibilité de cette mise à jour
comme histoire HMT.

On ne peut pas non plus déplacer seulement le registre et maintenir arbitrairement le
projecteur. La formule finie \eqref{lib01:P3} donne exactement
\[
 \|P_3S^3K\|^2=\frac{369035}{4}-\frac{281861\sqrt5}{10},
 \qquad
 \|P_3S^4K\|^2=\frac{1327717}{4}+\frac{694618\sqrt5}{5}.
\]
Elles sont distinctes de \eqref{lib01:normas-direccion} ; en particulier,
\([P_3,S^3]\) et \([P_3,S^4]\) ne s’annulent pas. En transportant
\(P_3\) par la même conjugaison, on conserve bien \(\eta_K\).

Le décodage est également covariant. Pour une permutation \(R\),
\[
 S'=RSR^{-1},\qquad M'=(R\oplus R)MR^{-1},\qquad
 M'R=(R\oplus R)M.
\]
Le réseau, la charge et les distances sont conservés. Dans le régime d’unicité,
la sortie du décodeur se transforme par \(R\) ; en dehors de celui-ci,
l’ensemble des minimiseurs est transporté, sans trancher arbitrairement
une égalité. Un repère orthogonal général exige de transporter également
\(\mathbb Z^{12}\) vers \(R\mathbb Z^{12}\) et la direction uniforme
vers \(R\mathbf1\). Arrondir dans l’ancien réseau après avoir changé le
système de coordonnées ne représente pas le même problème.

\begin{proposition}[Naturalité sous raffinement avec repères locaux]
\label{lib01:refinamiento}
Soit \(z\) un cylindre de successeurs \(z'\), de poids positifs compatibles
\(\sum_{\tau z'=z}\mu(z')=\mu(z)\) et de repères orthogonaux \(G_z\).
Le raffinement
\[
 \mathcal R_n(e_z\otimes x)=
 \sum_{\tau z'=z}\sqrt{\frac{\mu(z')}{\mu(z)}}
 e_{z'}\otimes G_{z'}G_z^*x
\]
est isométrique. Pour le champ de projecteurs
\(Q_z=G_zQ_KG_z^*\) et son opérateur par blocs \(Q_n\),
\[
 Q_{n+1}\mathcal R_n=\mathcal R_nQ_n,
 \qquad \mathcal R_n^*Q_{n+1}\mathcal R_n=Q_n.
\]
Les identités se composent entre niveaux et passent à la limite inductive.
\end{proposition}
\begin{proof}
Les successeurs ont des supports orthogonaux et les poids ont pour somme un ; les changements
de repère sont orthogonaux. Dans chaque successeur,
\(Q_{z'}G_{z'}G_z^*=G_{z'}G_z^*Q_z\), ce qui démontre le premier
carré. La multiplication par l’adjoint du raffinement donne le second.
La compatibilité par composition définit l’opérateur dans la limite des
inclusions isométriques. Le même argument vaut pour \(E_u\).
\end{proof}

L’énoncé conserve la représentation d’un préfixe sous des raffinements
compatibles. De nouveaux événements peuvent produire un autre registre ; on n’en déduit pas
une constance temporelle de \(K\). Les rangs des champs sur des fibres
tensorielles sont ceux de ces fibres, et non un décompte de dimensions physiques.

\begin{proposition}[Une charge de transport à action explicite]
\label{lib01:carga-noether}
Soient \(U_n:E_n\to E_{n+1}\) des transports inversibles. L’action
auxiliaire discrète à variables cotangentes est
\[
 \mathcal A(q,p)=\sum_n p_{n+1}^*(q_{n+1}-U_nq_n).
\]
Les variations à support fini donnent
\(q_{n+1}=U_nq_n\) et \(p_n=U_n^*p_{n+1}\). Si
\(G_0=I\), \(G_{n+1}=U_nG_n\), \(A_0=Q_K\) et
\(A_n=G_nA_0G_n^{-1}\), la quantité
\[
 J_n=p_n^*A_nq_n
\]
est indépendante de \(n\). L'action est invariante sous le groupe continu \(q_n\mapsto\exp(tA_n)q_n\), \(p_n\mapsto\exp(-tA_n^*)p_n\).
\end{proposition}
\begin{proof}
Les équations indiquées résultent de la variation de \(p_{n+1}\) et \(q_n\).
La définition de \(A_n\) implique
\(A_{n+1}U_n=U_nA_n\). Par conséquent,
\[
 p_{n+1}^*A_{n+1}q_{n+1}
 =p_{n+1}^*U_nA_nq_n=p_n^*A_nq_n.
\]
Le même entrelacement vaut pour les exponentielles et rend chaque terme de l'action invariant. Il n'est pas requis que \(A_n\) soit antisymétrique : la transformation du covecteur est la contragrédiente.
\end{proof}

Pour des transports orthogonaux, on peut choisir \(p_n=q_n\) ; alors
\(J_n=\|A_nq_n\|^2\), et le complément fournit la charge
restante. Dans une réalisation symplectique sans forme positive, on conserve
l’appariement écrit sans l’appeler énergie positive. Si l’on
archive les fibres au moyen d’isométries \(\mathcal V_n\), les
définitions
\[
 \widehat U_n=\mathcal V_{n+1}U_n\mathcal V_n^*,\qquad
 \widehat A_n=\mathcal V_nA_n\mathcal V_n^*,\qquad
 \widehat q_n=\mathcal V_nq_n,
\]
avec \(\widehat p_n(y)=p_n(\mathcal V_n^*y)\), conservent
l’entrelacement et la même charge dans les images. Cette composition
fournit une action et une symétrie continues concrètes pour une loi de
conservation ; elle n’attribue pas de courants de Noether automatiques à un groupe
fini et n’identifie pas l’action auxiliaire à toute action physique.

\subsection{Codage, identification opératorielle et stabilité de l'incidence}
\label{lib01:alcance}

La contribution réunie est une chaîne opératoire : le registre engendré possède un codage réversible ; son orbite est une référence complète ; le terminal et les compléments conservent cette référence et les observables transportés ; la séparation entière de deux mémoires permet de restituer l'incidence et la direction même lorsque le sélecteur réel a une marge nulle. Les preuves de linéarité, de projection et de réseaux sont utilisées comme langage de réalisation ultérieure, sans changer la production de l'objet.

Chaque hypothèse a un témoin qui empêche de l’élargir par inadvertance.
Une réponse isolée n’identifie pas un opérateur qui ne commute pas avec le
cycle ; seule la mémoire complète conserve automatiquement la matrice de Gram initiale ;
les termes croisés ne peuvent pas être omis ; la sélection absolue exige une
charge ; le projecteur requiert une direction non nulle ; les bornes de bruit
dépendent de la normalisation ; une transformation du système de coordonnées exige de transporter
la forme, le réseau et les lecteurs. Une perturbation de lecture ne devient pas par
décret une histoire admissible. Ces délimitations préservent les
fonctions de \(K\), sans le réduire à douze nombres déconnectés ni
promouvoir une propriété canonique au rang d’invariance universelle.
